% Generated by scripts/assemble.py; corrections belong in source/corrections.json.
\examhead{77}{吉林大学 2026 年数学分析试卷}{75}
% p076-b23; source PDF page 76
\par\Needspace{5\baselineskip}
\blockanchor{p076-b23}
\majorquestion{一}计算题.\par
% p076-b24; source PDF page 76
\blockanchor{p076-b24}
\question{1}$\lim_{n\to \infty}n^3\left[\sin \left(\tan \frac{1}{n}\right) - \sin \frac{1}{n}\right].$\par
% p076-b25; source PDF page 76
\blockanchor{p076-b25}
\question{2}$\lim_{n\to\infty}\frac{\sqrt[n]{(3n)!}}{n^{3}}.$\par
% p076-b26; source PDF page 76
\blockanchor{p076-b26}
\question{3}\displayformula{\lim _ {n \rightarrow \infty} \left[\left(\frac {\ln ((n + 1) !)}{n + 1} - \frac {\ln (n !)}{n}\right) \sum_ {k = 1} ^ {n} \sqrt [ k ]{k ^ {2} + \sin k} \right].}\par
% p077-b00; source PDF page 77
\blockanchor{p077-b00}
\question{4}$\lim_{x\to 0}\frac{(\sqrt{1 + x^2} - \sqrt{\cos x})\tan x^2}{\int_{x^2}^x\ln(1 + t^3)\mathrm{d}t}.$\par
% p077-b01; source PDF page 77
\blockanchor{p077-b01}
\question{5}$\lim_{n\to \infty}n\int_0^1 x^n\ln (1 + x)\mathrm{d}x.$\par
% p077-b02; source PDF page 77
\blockanchor{p077-b02}
\question{6}设 $u, v$ 是由 $\left\{ \begin{array}{l} u^2 - v = 3x + y \\ u - 2v^2 = x - 2y \end{array} \right.$ 确定的隐函数, 求 $\frac{\partial u}{\partial x}, \frac{\partial v}{\partial y}$.\par
% p077-b03; source PDF page 77
\blockanchor{p077-b03}
\question{7}计算 $\iiint_{D} (x^{2} + y^{2} + z^{2}) \, \mathrm{d}x \, \mathrm{d}y \, \mathrm{d}z$, 其中 $D$ 是两个球 $x^{2} + y^{2} + z^{2} \leq R^{2}$, $x^{2} + y^{2} + (z - R)^{2} \leq R^{2}$ 的公共部分.\par
% p077-b04; source PDF page 77
\blockanchor{p077-b04}
\question{8}计算第二型曲线积分 $\int_{L} \frac{y \, \mathrm{d}x + (1 - x) \, \mathrm{d}y}{(x - 1)^2 + y^2}$, 其中 $L$ 为圆周 $x^2 + y^2 = 4$, 方向取顺时针.\par
% p077-b05; source PDF page 77
\par\Needspace{5\baselineskip}
\blockanchor{p077-b05}
\majorquestion{二}当 $x, y, z$ 都大于 0 时, 求 $u = \ln x + 2\ln y + 3\ln z$ 在球面 $x^{2} + y^{2} + z^{2} = 6$ 上的最大值.\par
% p077-b06; source PDF page 77
\par\Needspace{5\baselineskip}
\blockanchor{p077-b06}
\majorquestion{三}证明函数项级数 $\sum_{n=1}^{\infty} \sin nx \ln \left(1 + \frac{1}{n}\right)$ 在 $(0,2\pi)$ 上逐点收敛但不一致收敛.\par
% p077-b07; source PDF page 77
\par\Needspace{5\baselineskip}
\blockanchor{p077-b07}
\majorquestion{四}求数项级数 $\sum_{n=1}^{\infty} \frac{1}{n(2n-1)} \left(-\frac{1}{4}\right)^n$ 的和.\par
% p077-b08; source PDF page 77
\par\Needspace{5\baselineskip}
\blockanchor{p077-b08}
\majorquestion{五}讨论参数 $p$ 对广义积分 $\int_0^{+\infty} \frac{\cos\left(x - \frac{1}{x}\right)}{x^p} \mathrm{d}x$ 的敛散性影响, 何时绝对收敛, 何时条件收敛, 何时发散?\par
% p077-b09; source PDF page 77
\par\Needspace{5\baselineskip}
\blockanchor{p077-b09}
\majorquestion{六}设 $f(x)$ 在 $[1,+\infty)$ 上有定义, 且存在常数 $L$,\par
% p077-b10; source PDF page 77
\par\Needspace{6\baselineskip}
\blockanchor{p077-b10}
使得对于任意 $x_1,x_2\in[1,+\infty)$, 都有\par
\nopagebreak[4]
% p077-b11; source PDF page 77
\blockanchor{p077-b11}
\displayformula{| f (x _ {1}) - f (x _ {2}) | \leq L | x _ {1} - x _ {2} |.}
% p077-b12; source PDF page 77
\blockanchor{p077-b12}
试证存在 $X \in [1, +\infty)$, 使得 $\frac{f(x)}{x + \ln\left(1 + \frac{1}{x}\right)}$ 在 $[X, +\infty)$ 上一致连续.\par
% p077-b13; source PDF page 77
\par\Needspace{5\baselineskip}
\par\Needspace{6\baselineskip}
\blockanchor{p077-b13}
\majorquestion{七}假设 $f(x)$ 在 $[0,2]$ 上四次连续可微, 设\par
\nopagebreak[4]
% p077-b14; source PDF page 77
\blockanchor{p077-b14}
\displayformula{g (x) = \frac {f (0)}{2} (x - 1) (x - 2) - f (1) x (x - 2) + \frac {f (2)}{2} x (x - 1).}
% p077-b15; source PDF page 77
\par\Needspace{6\baselineskip}
\blockanchor{p077-b15}
(1) 证明: 任给 $x \in [0,2]$, 都存在 $\xi \in (0,2)$, 使得\par
\nopagebreak[4]
% p077-b16; source PDF page 77
\blockanchor{p077-b16}
\displayformula{f (x) - g (x) = \frac {f ^ {\prime \prime \prime} (\xi)}{6} x (x - 1) (x - 2).}
% p077-b17; source PDF page 77
\blockanchor{p077-b17}
(2) 证明: 存在 $\eta \in (0,2)$, 使得 $\int_0^2 [f(x) - g(x)]\mathrm{d}x = -\frac{1}{90} f''''(\eta)$.\par
% p077-b18; source PDF page 77
\par\Needspace{6\baselineskip}
\blockanchor{p077-b18}
(3) 定义\par
\nopagebreak[4]
% p077-b19; source PDF page 77
\blockanchor{p077-b19}
\displayformula{\varepsilon_ {n} = \int_ {0} ^ {2} f (x) \mathrm{d} x - \frac {1}{3 n} \sum_ {i = 1} ^ {n} \left[ f \left(\frac {2 (i - 1)}{n}\right) + 4 f \left(\frac {2 i - 1}{n}\right) + f \left(\frac {2 i}{n}\right) \right]}
% p077-b20; source PDF page 77
\blockanchor{p077-b20}
证明： $\lim_{n\to \infty}n^4\varepsilon_n = -\frac{1}{180} [f'''(2) - f'''(0)].$\par
